%
%! regex-fix=true
%! dir=rataflechera

\songtitle{Throne of Smoke}
\tag*{2026}\tag{irish-folk-rock}\tag{original-lyrics}\tag{african-empire}\tag{monukae}\tag{abulon}\tag{gladiel}
\begin{notes}
Written by Carlos Thompson -- Gladiel's full arc within the \emph{Monukae} saga, from the Irish pirate king's daughter captured after her village was razed, through her mentorship under Abulon and marriage to Mombuto, to her rise as the secret queen of the sea. Expands on the backstory introduced in \emph{The Pirate's Ledger}. Retrieved from rataflechera's Suno page (V6, 16 September 2026).
\end{notes}
\begin{style}
Irish epic ballad, cinematic folk-rock with tape saturation, plate reverb, and a stereo-wide chorus; acoustic strumming, overdriven electric accents, frame drum pulse, West African percussion, ngoni plucks, kora filigree, brass stabs, and driving kick drum; brisk 132 BPM with a driving pulse and chant refrain; female lead in rapid-fire narrative delivery with layered unison vocals.
\end{style}
\begin{lyrics}
[Verse 1]
When she saw him, she was intrigued
a good specimen, elegant clothes torn in the battle
Gladiel fantasized with him that night
but soon forgot about him,
just one more prisoner of his father
a ransom from his wealthy family
or a slave for the highest bidder
not the kind of man that would stay in Munster
not the kind of man that Gladiel would get

[Pre-Chorus]
The second time she saw him she hated him
for his sake her village was razed and burning
for his sake her mother, aunts and uncles were slayed
for his sake she and her brother were captured

[Chorus]
Gladiel didn't choose her life
the daughter of the pirate
the wife of whoever her father would pick
but the proud free bird until that day would come
Gladiel didn't choose this other life
the prisoner, the mentee, the companion
the woman behind a throne in the distance
the woman with a throne of smoke by her own

[Verse 2]
Gladiel found herself in a golden cage
with open doors but miles inland in the Sahel
and a long sea voyage from a home that wasn't more
with daily dinner with the butchers of her family
she saw her father brought in chains moons after
cast away in some dungeon she was not let visit
then her brother came with their father's last words
erratic words before his own son, her brother, slayed him.

[Pre-Chorus]
a hunger strike, dying was the life Gladiel chose
then Abulon, her captor, came and talk to her
not as a child, not as a damsel, but as an equal
the father of the man that caused her fate,
offer to mentor her, not as a prisoner, not as a daughter
but as an equal partner.

[Chorus]
Gladiel didn't choose her life
the daughter of the pirate
the wife of whoever her father would pick
but the proud free bird until that day would come
Gladiel didn't choose this other life
the prisoner, the mentee, the companion
the woman behind a throne in the distance
the woman with a throne of smoke by her own

[Bridge]
the heart that once belong to her father
was now opening for Abulon
the golden prison was becoming her domain
the specimen that caused all trouble,
Mombutu, Abulon's son, a fun peer
and Marco, the Venetian mercenary,
the one waking up the woman in this girl.

[Chorus]
Gladiel didn't choose this other life
the prisoner, the mentee, the companion
the woman behind a throne in the distance
the woman with a throne of smoke by her own

[Verse 3]
Mombuto went to his grandparents in Peru and Marco with him
Mombuto came back married, Marco didn't come back
Gladiel didn't know who of these men hurt her more
but at the side of Abulon she would endure
Then Abulon told her she was ready, her training complete
An advisor to Mombuto to serve both in the Frankish court
A practical choice but also mutual respect, under Christian Rite
Mombuto and Gladiel became in Europe husband and wife

[Pre-Chorus]
Back in the Sahel, Numak, the other Mombutu's wife
just made it clear, her children were the heirs of Monukae
the most fertile land of the Sahel, the estate of Abulon
Gladiel complied, she was not made for the land

[Chorus]
Gladiel didn't choose her life
the daughter of the pirate
the wife of whoever her father would pick
but the proud free bird until that day would come
Gladiel didn't choose this other life
the prisoner, the mentee, the companion
the woman behind a throne in the distance
the woman with a throne of smoke by her own

[Outro]
Abulon died in his final battle, in the seas of the Romans
Gladiel cried for Abulon what she didn't cry for her father
The emperor ensured Mombuto wouldn't concentrate the power of Abulon
Mombuto left Monukae to Numak and his daughter, cared by his brothers
he left the fleet under Gladiel, while playing diplomat for the emperor.
the merchant navy was surrendered to the emperor
and a new secret merchant navy was left to Gladiel to rebuild
a throne of smoke but a real power
the daughter of the pirate, now the queen of the sea.
\end{lyrics}

\songtitle{Laoi Ghláidial}
\tag*{2026}\tag{sean-nos}\tag{irish-lyrics}\tag{african-empire}\tag{monukae}\tag{abulon}\tag{gladiel}\tag{throne-of-smoke}
\begin{notes}
Irish-language (Gaeilge) translation and adaptation by Claude of \emph{Throne of Smoke}, carrying Gladiel's full arc within the \emph{Monukae} saga into a sean-nós duet arrangement. Retrieved from rataflechera's Suno page (V6, 16 September 2026).
\end{notes}
\begin{style}
Irish epic ballad, female mezzo-soprano and male tenor duet with rapid-fire delivery, sean-nós ornamentation and layered choir vocals; 88 bodhrán pulse, half-time chorus; uilleann pipes drone, low whistle countermelody, West African djembe, kora arpeggios, talking drum accents, frame drum rolls, brass fanfare stabs; plate reverb, tape saturation, sidechain pumping, tragic resolve
\end{style}
\begin{lyrics}
[Verse 1]
An chéad uair dá bhfaca sí é,
a chulaith uasal stróicthe ón gcath;
fear breá, is rinne sí brionglóid
an oíche sin faoi, gan rath.

Ach rinne sí dearmad go luath air —
príosúnach eile ag a hathair,
airgead fuascailte óna mhuintir shaibhir,
nó daor le díol ar an margadh.

Ní fear é a d'fhanfadh sa Mhumhain,
ní fear é a gheobhadh Gláidial.

[Pre-Chorus]
An dara huair a chonaic sí é, ba ghráin léi é.
Dá bharr bhí a baile ina lasair is ina luaith,
dá bharr do maraíodh a máthair is a muintir,
dá bharr do gabhadh í féin is a deartháir.

[Chorus]
Níor thogh Gláidial a saol —
iníon an fhoghlaí mhara,
bean ag cibé fear a roghnódh a hathair,
ach éan saor uaibhreach go dtiocfadh an lá.

Níor thogh Gláidial an saol eile seo —
an cime, an dalta, an compánach,
an bhean taobh thiar de ríchathaoir i gcéin,
an bhean le ríchathaoir deataigh dá cuid féin.

[Verse 2]
Fuair Gláidial í féin i gcás óir,
na doirse ar oscailt, ach na mílte míle ón muir,
i dtír loiscthe an deiscirt, is turas fada farraige
ó bhaile nach raibh ann níos mó.

Dinnéar gach lá le buíon a mharaigh a muintir.
Chonaic sí a hathair i slabhraí, míonna ina dhiaidh sin,
caite i bpríosún nach ligtí di cuairt a thabhairt air.
Tháinig a deartháir le focail dheireanacha a n-athar —
focail mhearaí os comhair a mhic féin,
an mac a mharaigh é.

[Pre-Chorus]
Troscadh chun báis — ba é sin an saol a thogh Gláidial.
Ansin tháinig Abulón, a gabhálaí, chun cainte léi:
ní mar leanbh, ní mar mhaighdean, ach mar chomhionann.
Athair an fhir ba chúis lena cinniúint
ag tairiscint teagasc a thabhairt di —
ní mar chime, ní mar iníon, ach mar pháirtí.

[Chorus]
Níor thogh Gláidial a saol —
iníon an fhoghlaí mhara,
bean ag cibé fear a roghnódh a hathair,
ach éan saor uaibhreach go dtiocfadh an lá.

Níor thogh Gláidial an saol eile seo —
an cime, an dalta, an compánach,
an bhean taobh thiar de ríchathaoir i gcéin,
an bhean le ríchathaoir deataigh dá cuid féin.

[Bridge]
An croí ba le hathair uair amháin,
tá sé ag oscailt anois d'Abulón.
Tá an príosún óir ag éirí ina ríocht.

An fear breá ba chúis leis an mbuairt ar fad —
Mombutú, mac Abulóin, compánach suairc;
agus Marcó, an t-amhas Veinéiseach,
an té a dhúisigh an bhean sa chailín seo.

[Chorus — half]

Níor thogh Gláidial an saol eile seo —
an cime, an dalta, an compánach,
an bhean taobh thiar de ríchathaoir i gcéin,
an bhean le ríchathaoir deataigh dá cuid féin.

[Verse 3]
Chuaigh Mombutú go dtí a shinsir i bPeiriú, is Marcó leis.
Tháinig Mombutú ar ais pósta; níor tháinig Marcó ar ais choíche.
Níorbh eol do Ghláidial cé acu fear ba mhó a ghoill uirthi,
ach le hais Abulóin d'fhulaingeodh sí.

Ansin dúirt Abulón léi go raibh sí ullamh, a hoiliúint slán:
comhairleach do Mhombutú, an bheirt acu i gcúirt na bhFranc.
Rogha phraiticiúil, ach urraim fhrithpháirteach leis —
faoi dheasghnáth Chríostaí, san Eoraip i gcéin,
rinneadh fear céile is bean chéile díobh.

[Pre-Chorus]
Thiar i dtír loiscthe an deiscirt, Númac, an bhean eile ag Mombutú,
chuir sí in iúl go soiléir: ba iad a clann-sean oidhrí Mhonúcae,
an talamh ba thorthúla sa deisceart, eastát Abulóin.
Ghéill Gláidial. Níor tógadh ise don talamh.

[Chorus]
Níor thogh Gláidial a saol —
iníon an fhoghlaí mhara,
bean ag cibé fear a roghnódh a hathair,
ach éan saor uaibhreach go dtiocfadh an lá.

Níor thogh Gláidial an saol eile seo —
an cime, an dalta, an compánach,
an bhean taobh thiar de ríchathaoir i gcéin,
an bhean le ríchathaoir deataigh dá cuid féin.

[Outro]
Fuair Abulón bás ina chath deireanach, i bhfarraigí na Róimhe.
Chaoin Gláidial é mar nár chaoin sí a hathair riamh.

Chinntigh an tImpire nach mbeadh cumhacht Abulóin i lámha aon fhir amháin:
d'fhág Mombutú Monúcae ag Númac is ag a iníon,
faoi chúram a dheartháireacha;
d'fhág sé an cabhlach faoi Ghláidial,
is é féin ag imirt an taidhleora don Impire.

Tugadh an loingeas trádála suas don Impire —
is fágadh loingeas rúnda nua ag Gláidial le tógáil as an nua.

Ríchathaoir deataigh, ach fíorchumhacht.
Iníon an fhoghlaí mhara —
banríon na mara anois.
\end{lyrics}

%! dir=gabisson

\songtitle{It's Still With Me}
\tag*{2026}\tag{chamber-pop}\tag{original-lyrics}
\begin{notes}
Written by Carlos Thompson -- faith and apostasy argued point by point, what survives of Christian ethics after leaving the belief behind, and a pointed critique of Christians who wield the label as tribal identity rather than practice. Retrieved from Gabriel's Suno page (V6-MINI, 15 September 2026).
\end{notes}
\begin{style}
Chamber pop with country-rock verses: fingerpicked acoustic guitar, pedal steel swells, electric piano pads, brushed floor toms, group claps, and marching snare at 92 BPM; male baritone with conversational spoken intro, restrained verses, and a hymnlike gospel-choir chorus, shaped by plate reverb, tape saturation, and a clean hi-fi mix.
\end{style}
\begin{lyrics}
[Verse 1]
I can't do religions, I'm no longer in the church
It is not about your god, I never fought against it
It is not about the flock, you're just people like me
I just can't pretend obvious fables are factuality

I was raised Christian and that still means something to me
The values I praise come from the teachings of the Christ
They taught me there was no humanism without Christianity
I can no longer believe in the creed, I still hold humanity

[Pre-Chorus]
My church taught me to love my neighbor
To see the divine in my fellow people
And that's something that transcended
After my apostasy, it's still with me

[Chorus]
How can you call yourself a Christian
If you take the label as a badge
A tribe behind a cross and a Bible
Not the teachings of a loving god

How can you call yourself a Christian
To spread hate instead of love
To find ways to judge the other
And bear the beam in your eye with pride

[Verse 2]
You vote against social support networks and call it fine
No government should force your charity, that's with you and god
When Jesus mandated to feed the poor, what was in his words?
Was that a test of your piety? Or a call not to leave the poor behind?

You demand infidels to remain outside your borders
That's not what the old book taught, nor the words of the Christ
You're caring more for your Christian label, for the forms
For the purity of your tribe than what the faith is for

[Pre-Chorus]
My church taught me to take care of others
To fight for a world with less injustice
The suffering of my neighbor is my cross to bear
After my apostasy, it's still with me

[Chorus]
How can you call yourself a Christian
If you take the label as a badge
A tribe behind a cross and a Bible
Not the teachings of a loving god

How can you call yourself a Christian
To spread hate instead of love
To find ways to judge the other
And bear the beam in your eye with pride

[Bridge]
Did they teach you not to care for others?
Call those who fail moral failures?
That the wealthy deserve tax cuts
Because that's the will of your god?

Love god above everything
And thy neighbor as thou lovest thyself
Love thy enemy, offer the other cheek
But only when it doesn't bother thee

Just tell yourself that further thinking
Is the deal of the devil
Be servile to the powerful
Let the poor rot in mud

[Final Chorus]
How can you call yourself a Christian
If you take the label as a badge
A tribe behind a cross and a Bible
Not the teachings of a loving god

How can you call yourself a Christian
To spread hate instead of love
To find ways to judge the other
And bear the beam in your eye with pride

I don't call myself a Christian
I just can't believe in the fables
I still think the fables bear wisdom
After my apostasy,
It's still with me
\end{lyrics}

\songtitle{Manet Id}
\tag*{2026}\tag{gregorian-chant}\tag{latin-lyrics}\tag{its-still-with-me}
\begin{notes}
Latin-language translation and adaptation by Claude of \emph{It's Still With Me}, carrying the same faith-and-apostasy argument into an a cappella Gregorian-chant setting. Retrieved from Gabriel's Suno page (V6-MINI, 15 September 2026).
\end{notes}
\begin{style}
Gregorian chant, unison male choir and sustained female choir drone, entirely a cappella with no instruments, single-microphone ensemble capture, austere processional march with measured modal phrasing, stern sacred choral delivery, restrained room resonance and natural ensemble blend
\end{style}
\begin{lyrics}
[Versus I]
Iam fidem non teneo,
Ecclesiae non adsum;
Non de Deo tuo agitur,
Numquam illi adversatus sum.

Non de grege est sermo meus,
Homo es sicut ego sum;
Fabulas simulare nequeo
Claras, tamquam sint verum.

Educatus sum Christianus,
Hoc adhuc mihi valet;
Virtutes quas laudo veniunt
Ex doctrina quam Christus dedit.

Docuerunt me sine Christo
Humanitatem non stare;
Credere iam symbolo non possum,
Sed humanitatem teneo care.

[Praeludium chori]
Ecclesia amare proximum
Docuit me olim,
Divinum in fratre cernere —
Manet id post apostasiam.

[Chorus]
Quomodo te dicis Christianum,
Si nomen sumis ut insigne,
Tribum post crucem et biblia,
Non doctrinam Dei amantis?

Quomodo te dicis Christianum,
Odium spargens pro amore,
Alium semper iudicando,
Trabem in oculo tuo superbe ferens?

[Versus II]
Contra auxilium sociale votas,
Et dicis rem esse bonam;
Nulla res publica cogat pietatem tuam,
Id tuum est cum Deo tantum.

Cum Iesus iussit pascere egenos,
Quid verba illius volebant?
Fuitne probatio pietatis tuae?
An vox ne pauperes relinquerentur?

Infideles extra fines tuos
Manere iubes semper;
Non hoc liber vetus docuit,
Nec verba Christi dixerunt.

Magis curas de nomine Christiano,
De formis, de specie,
De tribus tuae puritate
Quam de causa fidei.

[Praeludium chori II]
Ecclesia curam aliorum
Docuit me habere,
Pro mundo minus iniquo pugnare —
Dolor proximi crux mea est.
Post apostasiam meam
Hoc mecum permanet.

[Chorus]
Quomodo te dicis Christianum,
Si nomen sumis ut insigne,
Tribum post crucem et biblia,
Non doctrinam Dei amantis?

Quomodo te dicis Christianum,
Odium spargens pro amore,
Alium semper iudicando,
Trabem in oculo tuo superbe ferens?

[Transitus]
Docueruntne te non curare alios?
Vocare peccatores eos qui cadunt?
Divites merere tributa minora,
Quia haec voluntas Dei tui est?

Dilige Deum super omnia,
Et proximum sicut te ipsum;
Dilige inimicum, praebe alteram genam,
Sed tantum cum id te non gravat.

Dic tibi ipsi: ultra cogitare
Diaboli esse negotium;
Esto servilis potentibus,
Sine pauperes in luto putrescere.

[Chorus finalis]
Quomodo te dicis Christianum,
Si nomen sumis ut insigne,
Tribum post crucem et biblia,
Non doctrinam Dei amantis?

Quomodo te dicis Christianum,
Odium spargens pro amore,
Alium semper iudicando,
Trabem in oculo tuo superbe ferens?

Non me dico Christianum,
Fabulis credere non possum;
Adhuc puto fabulas sapientiam ferre —
Post apostasiam meam,
Mecum adhuc manet.
\end{lyrics}

%! dir=wyomee

\songtitle{لایە‌لایەی جان‌ئالان — Laye-layey Jan-Alan}
\tag*{2026}\tag{kurdish-lullaby}\tag{kurdish-lyrics}\tag{thomas-entourage}\tag{zahra-moradi}\tag{mihai-rusu}\tag{jan-alan}
\begin{notes}
From the \emph{Thomas' Entourage} narrative universe by Carlos Thompson -- a Kurdish lullaby for Jan-Alan, Zahra Moradi and Mihai Rusu's son, tracing his family's geography from Kraków to Sanandaj to his grandparents' hills near Ismail. Kurdish translation by Claude of English lyrics written by SUNO. Retrieved from Wyomee Dank's Suno page (V6-MINI, 16 September 2026).
\end{notes}
\begin{style}
Kurdish lullaby, gentle female mezzo-soprano, slow rocking 6/8 pulse, intimate acoustic arrangement with santur, soft tembûr, warm frame drum brushed very lightly, delicate piano, spacious but close vocal recording, simple memorable melody, soothing bedtime dynamics, subtle cradle-like swells between sections
\end{style}
\begin{lyrics}
[بەشی یەکەم / Verse 1]
جان‌ئالان، چاوەکانت دابخە، بچکۆلەکەم
ئاسمانی هێمنی کراکۆڤ تەماشات دەکات
هەورێکی زیوین بە پەنجەرەدا تێدەپەڕێت
هەر هەناسەیەک ئاشتیت بۆ دەهێنێت

[کۆرەس / Chorus]
بنوە، جان‌ئالان، بە نەرمی بنوە
با خەونی خۆش لە نێوچاوانت کۆببنەوە
لە ئاسمانی کراکۆڤەوە بۆ ڕێگا دوورەکان
خۆشەویستی لەگەڵتدا دەڕوات

[بەشی دووەم / Verse 2]
تۆ گردەکانی ئیزمائیلت بینی
لەوێ کەسوکارت گەرم لە باوەشیان گرتیت
پێکەنینیان بە نەرمی بۆ ماڵەوە هێنایتەوە
گەنجینەیەکی بچووکی ڕووناک لە گوێتدا

[کۆرەس / Chorus]
بنوە، جان‌ئالان، بە نەرمی بنوە
با خەونی خۆش لە نێوچاوانت کۆببنەوە
لە ئیزمائیلەوە بۆ سنە
خۆشەویستانت هەر هەنگاوێکت پیرۆز دەکەن

[پردە / Bridge]
بای سنە ناوە شیرینەکەت دەزانێت
بە نەرمی بە شەودا دەیچڕێت
دەستی داپیرە و دوعای باپیرە
چرایەکی بچووک بە ڕووناکی دەهێڵنەوە

[کۆتایی / Outro]
جان‌ئالان، لەژێر ئاسمانی پۆڵەندا پشوو بدە
کراکۆڤ هێمنە، مانگ بەرزە
خەونە بچووکەکانت بۆ هەر شوێنێک بڕۆن
خۆشەویستی چاوەڕێت دەکات، بۆ ماڵەوە ڕێنماییت دەکات
\end{lyrics}

%! dir=gabisson

\songtitle{Twenty Twelve Scripts}
\tag*{2012}\tag{kuduro}\tag{original-source}
\begin{notes}
Lyrics by Suno based on a summary of Carlos Thompson's 2012 blog production, written under his own name -- essentially a year of post titles.
\end{notes}
\begin{style}
Angolan kuduro, fast syncopated electronic percussion, clipped 4/4 kick pattern, rubbery bass synth, bright metallic stabs, call-and-response crowd chants, tense playful political energy, English male lead vocal with rhythmic half-spoken delivery, sudden stop-start breaks, compact verse-to-hook arrangement, dry punchy club mix with sharply separated drums
\end{style}
\begin{lyrics}
[Intro]
Twenty twelve scripts
Keep the page under lock
Too soon for the chapter
But the clock won't stop

[Verse 1]
January, steam in the kitchen
A sealed-up word with a loud conviction
February, doubt on the screen
Everybody certain, nobody clean
March brought talent left in the rain
A blocked-up road and a borrowed lane
I crossed your line, you crossed mine
Red signs blinking, still we align

[Pre-Chorus]
Read the margin, count the gaps
History hides inside the scraps
Say less, move quick, hold that thought
What isn't published still gets fought

[Chorus]
Twenty twelve scripts, keep it moving
No final page, but the truth keeps proving
Twenty twelve scripts, don't ask why
Ink in the drawer, fire in the sky
We don't need the whole story tonight
We dance through the missing lines
Twenty twelve scripts, hold that beat
A secret year with a public street

[Verse 2]
April, legal thoughts in a crowded room
May on campus, making speculation bloom
Peace wears left, peace wears right
Same old question under different lights
A pirate sea between two shores
Promises knocking on guarded doors
June asked who guarantees the deal
While the city rolled on four small wheels

[Bridge]
Move, move, move through the question
Stop, stop, stop the possession
Freedom isn't one clean road
Carry the weight, unload the code

[Chorus]
Twenty twelve scripts, keep it moving
No final page, but the truth keeps proving
Twenty twelve scripts, don't ask why
Ink in the drawer, fire in the sky
We don't need the whole story tonight
We dance through the missing lines
Twenty twelve scripts, hold that beat
A secret year with a public street

[Verse 3]
August wind through a restless town
Peace came up, then the walls came down
Liberty arguing with liberty's face
Every promise running its own race
December turns the wheel once more
Old leaves waiting by the door
Cycles, recycle, start again
The year closes softly, but it won't end

[Outro]
Too soon to publish, too loud to erase
A hidden chapter still leaves a trace
Twenty twelve scripts
Keep the page under lock
When the seal breaks open
We will hear the clock
\end{lyrics}

%! dir=gabisson

\songtitle{Twenty Thirteen Scripts}
\tag*{2013}\tag{forro}\tag{original-source}
\begin{notes}
Lyrics by Suno based on a summary of Carlos Thompson's 2013 blog production, written under his own name -- essentially a year of post titles.
\end{notes}
\begin{style}
Brazilian forró pé-de-serra with English lyrics, lively accordion lead, zabumba and triangle groove, nylon-string guitar, warm male vocal, playful scholarly storytelling, bright danceable verses, brief reflective breakdowns between sections, organic room sound, festive Latin alternative-rock lift in the final chorus, clear melodic hooks and a joyful Christmas coda
\end{style}
\begin{lyrics}
[Intro]
Twenty thirteen, turn the page, come on
A little site, a restless age, come on

[Verse 1]
February asks what can be believed
A secondary world that won't be retrieved
March writes down what I want to become
A soul-shaped guess beneath the sun
April brings two tongues to the table
One says assume, one says know if you're able
An agnostic self with a notebook spine
Learning where the question crosses the line

[Chorus]
Twenty thirteen, swing the door wide
Ideas in boots, moving side to side
Chlewey scripts, still too soon to print
But every unfinished thought leaves a glint
Twenty thirteen, don't wait for proof
Dance with the doubt beneath the roof
If history's breathing, let it begin
A little more honest, a little less thin

[Verse 2]
May takes a journey past the old refrain
June asks whether logic can carry the rain
Materialism tired, but the mind stays awake
An apathetic god is a difficult shape
August says, "Pass the testimony—stop or go?"
A witness at the crossing with a question in tow
September presupposes morals in the dust
What do we owe each other before we trust?

[Chorus]
Twenty thirteen, swing the door wide
Ideas in boots, moving side to side
Chlewey scripts, still too soon to print
But every unfinished thought leaves a glint
Twenty thirteen, don't wait for proof
Dance with the doubt beneath the roof
If history's breathing, let it begin
A little more honest, a little less thin

[Bridge]
October turns the calendar over slowly
November brings the White Mister, speaking boldly
"What game is she playing?" the indignant tweeter cries
While a papaya-rock answer rolls under the skies
No final doctrine, no polished decree
Just notes from a life trying hard to see

[Verse 3]
The blog keeps walking through the questions at night
Between Spanish phrases and a borrowed light
What is a self when the languages divide?
What stays real when the second world opens wide?
The scripts wait quietly, not ready for release
But the waiting itself has a rhythm and peace
And every small heading, each unfinished design
Becomes a road marker in the history line

[Final Chorus]
Twenty thirteen, swing the door wide
Ideas in boots, moving side to side
Chlewey scripts, let the old pages sing
Too soon for publishing, not too soon for spring
Twenty thirteen, raise your glass high
For the questions that refuse to die
From February rain to December cheer
The unfinished future is already here

[Outro]
Christmas on the papaya, bells in the street
Rock in Spanish with a forró beat
Twenty thirteen, the page turns bright
See you in the next unfinished night
\end{lyrics}

%! dir=gabisson

\songtitle{Twenty Fourteen Scripts}
\tag*{2014}\tag{dancehall}\tag{original-source}
\begin{notes}
Lyrics by Suno based on a summary of Carlos Thompson's 2014 blog production, written under his own name -- essentially a year of post titles.
\end{notes}
\begin{style}
Dancehall fused with bright Jamaican ska, brisk offbeat guitar chops, springy bass, crisp rimshot and hand percussion, occasional brass stabs, playful conversational English lead vocal, chant-ready chorus, compact verse-to-chorus arrangement, warm punchy modern mix with lively horn accents
\end{style}
\begin{lyrics}
[Intro]
Twenty fourteen scripts, page still turning
Chlewey on the typewriter, questions burning
Not for release, the envelope stays sealed
But the year left footprints in the field

[Verse 1]
January: knowing, believing, two roads in the dust
March: divine karma, why existence earns trust
April counting figures, making numbers behave
May four more years, promises we save

[Pre-Chorus]
Read the headings, feel the change
Every month rearranged
Some things written, some concealed
History moves before it's revealed

[Chorus]
Twenty fourteen scripts, keep the paper dry
Pages in a locked room, waiting for the sky
Twenty fourteen scripts, let the riddles run
What was too soon to publish still remembers everyone

[Verse 2]
July: sacrificing children, a warning in the frame
September: prohibiting water, who gets to make the claim?
Challenges in the margin, ink beneath the line
October asks the question of a people and a sign

[Pre-Chorus]
Read the headings, feel the change
Every month rearranged
Some things written, some concealed
History moves before it's revealed

[Chorus]
Twenty fourteen scripts, keep the paper dry
Pages in a locked room, waiting for the sky
Twenty fourteen scripts, let the riddles run
What was too soon to publish still remembers everyone

[Bridge]
No spoiler, no verdict, no borrowed crown
Just a blog in the archive while the years roll down
If the future opens what the present hid
We'll meet those old questions like we always did

[Dancehall Chant]
January, March, April, May
July, September, October—say
Write it down but hold it close
A sealed-up year becomes a ghost

[Final Chorus]
Twenty fourteen scripts, keep the paper dry
Pages in a locked room, waiting for the sky
Twenty fourteen scripts, let the riddles run
What was too soon to publish still remembers everyone

[Outro]
Chlewey's page, the bookmark stays
2014 in a stack of days
No reveal tonight, just a rhythm and a trace
History keeps dancing till it finds its place
\end{lyrics}

%! dir=gabisson

\songtitle{Twenty Fifteen Scripts}
\tag*{2015}\tag{porro-cumbia}\tag{original-source}
\begin{notes}
Lyrics by Suno based on a summary of Carlos Thompson's 2015 blog production, written under his own name -- essentially a year of post titles.
\end{notes}
\begin{style}
Porro-cumbia, lively Colombian Caribbean groove, syncopated hand percussion, guacharaca scrape, alegre and tambora drums, bright brass melodies, accordion flourishes, buoyant bass, playful English lead vocal, communal call-and-response chorus, crisp festive mix, verses with nimble storytelling and choruses opening into brass-led celebration
\end{style}
\begin{lyrics}
[Intro]
Open the folder, turn the page,
A little old world in a digital age.

[Verse 1]
January came with a welcome sign,
Peace in the title and questions in line.
Four small windows, a room full of news,
A joke in French and a point of view.
February brought four names to the door,
José, Jaime, Jacobo, Joaquim, and more.
Our indignation wore a borrowed coat,
Then sailed downriver in a paper boat.

[Pre-Chorus]
Some things were written, some stayed sealed,
Some truths waited for time to yield.

[Chorus]
Twenty fifteen scripts, everybody turn the page,
Notes from a restless year, dancing through the digital age.
From the morning post to the river mist,
What we could publish made the list.
Twenty fifteen scripts, let the old ideas persist,
A little history with a Caribbean twist.

[Verse 2]
April said indignation everywhere,
War in the distance, smoke in the air.
July found a ride back home,
Uberreflexiones on the road alone.
Then Ugenito entered, saga in hand,
A tiny hero with a stubborn plan.
He crossed each paragraph, bright-eyed and keen,
Looking for the line between what's said and seen.

[Pre-Chorus]
Some things were written, some stayed sealed,
Some truths waited for time to yield.

[Chorus]
Twenty fifteen scripts, everybody turn the page,
Notes from a restless year, dancing through the digital age.
From the morning post to the river mist,
What we could publish made the list.
Twenty fifteen scripts, let the old ideas persist,
A little history with a Caribbean twist.

[Bridge]
August smelled of river water,
A fish watched the world like a quiet daughter.
Alternative cures and cities unknown,
Ten far-off places, still not home.
Pain for a country, love underneath,
A question kept warm behind the teeth.
And the deepest chapter, held from view,
Was still too soon for me and you.

[Break]
Hey—save that thought!
Hey—leave that line!
The future may open it in time!

[Final Chorus]
Twenty fifteen scripts, everybody turn the page,
Notes from a restless year, dancing through the digital age.
From the morning post to the river mist,
What we could publish made the list.
Twenty fifteen scripts, let the old ideas persist,
A little history with a Caribbean twist.

[Outro]
Close the folder, keep the key,
Some stories wait to become history.
Twenty fifteen, still humming in the air—
A page not published, but already there.
\end{lyrics}

%! dir=gabisson

\songtitle{Let the Border Breathe Again}
\tag*{2015}\tag{tropical-pop}\tag{original-source}
\begin{notes}
Lyrics by Suno based on a Claude summary of Carlos Thompson's 2015 blog post \emph{Dolor de patria} -- a critique of nationalism and border politics, migration, xenophobia, and the manufactured outrage that turns human displacement into a political weapon on both sides. Retrieved from Gabriel's Suno page (V6-MINI, 17 September 2026).
\end{notes}
\begin{style}
Tropical pop, candy-bright synths, buoyant dembow-inspired groove, crisp hand percussion, elastic bass, acoustic guitar accents, playful melodic female lead, stacked sugary harmonies, polished radio mix, verses delivered conversationally, pre-chorus building tension, wide singalong chorus, brief reflective bridge, bittersweet but energetic
\end{style}
\begin{lyrics}
[Verse 1]
At the checkpoint, names become numbers
A suitcase pressed against a wall
Some came for bread, some fled the thunder
Some had papers, some had none at all

[Pre-Chorus]
A flag can flare in every window
A headline can decide who's right
But tell me who counts every doorway
When home disappears overnight?

[Chorus]
Don't turn pain into a parade
Don't sell wounded pride by the page
If one child crosses and one child's displaced
Both deserve a hand, not a cage
We can love where our mothers were raised
Without making a stranger the blame
No more borrowed fury, no more patriotic game
Let the border breathe again

[Verse 2]
Cheap fuel, empty shelves, long lines waiting
A neighbor sees an outsider's gain
When fear is dressed up as protection
Every difference starts to look like shame

[Pre-Chorus]
The law may draw a clean black border
Life keeps writing in the gray
And leaders wave the loudest colors
When an election's on the way

[Chorus]
Don't turn pain into a parade
Don't sell wounded pride by the page
If one child crosses and one child's displaced
Both deserve a hand, not a cage
We can love where our mothers were raised
Without making a stranger the blame
No more borrowed fury, no more patriotic game
Let the border breathe again

[Bridge]
Three thousand homes can vanish quietly
While cameras chase the next decree
A handshake takes more courage than a battle cry
And mercy is not betrayal of the country

[Final Chorus]
Don't turn pain into a parade
Don't sell wounded pride by the page
If one child crosses and one child's displaced
Both deserve a hand, not a cage
We can love where our mothers were raised
And still leave room for someone else's name
No more borrowed fury, no more patriotic game
Let the border breathe again

[Outro]
Let the border breathe again
Let the border breathe again
\end{lyrics}

\songtitle{No me vendas dolor de patria}
\tag*{2015}\tag{tropical-pop}\tag{spanish-lyrics}\tag{original-source}\tag{let-the-border-breathe-again}
\begin{notes}
Lyrics by Suno based on a Claude summary of Carlos Thompson's 2015 blog post \emph{Dolor de patria}. Spanish-language companion to \emph{Let the Border Breathe Again}, generated in the same session -- localizes the argument to the Colombia--Venezuela border, naming the Corinto eviction of a Nasa family and four Colombian presidents by name. Retrieved from Gabriel's Suno page (V6-MINI, 17 September 2026).
\end{notes}
\begin{style}
Tropical pop with candy-pop hooks, upbeat dembow-inspired groove, bright marimba plucks, clean electric guitar, bubbly synth bass, handclaps, crisp modern radio production; conversational Spanish lead vocal with playful verses, tight melodic chorus, brief half-time bridge, energetic final refrain
\end{style}
\begin{lyrics}
[Intro]
No todo grito es verdad,
no toda bandera sabe cuidar.

[Verse 1]
Primero, pon los nombres sobre la mesa,
no todo colombiano es guerra ni amenaza.
Hay grupos ilegales, sí, cruzando la frontera,
pero un sello en un discurso no explica la historia entera.

Segundo, una expulsión puede tener ley,
los gobiernos deportan, aquí también lo sé.
En Corinto cayó una casa nasa aquel año,
y el silencio pesó más que cualquier titular extraño.

[Pre-Chorus]
Pero entre tratados viejos y cédulas dudosas,
la vida no cabe en frases luminosas.
Libre tránsito, voto comprado, papeles sin final,
la frontera es una herida con memoria legal.

[Chorus]
No me vendas dolor de patria,
si eliges qué dolor mirar.
Tres mil dentro de la casa,
¿quién los vino a reclamar?
No me vendas orgullo herido,
con la rabia en oferta especial.
Si la bandera tapa al vecino,
¿de qué sirve gritar nacional?

[Verse 2]
País A es pobre, país B tiene más,
X busca trabajo, Y teme perder su lugar.
Subsidios, gasolina, mercado y contrabando,
el que llega se vuelve invasor en el relato.

Pero X también puede huir de una finca,
de una guerrilla, de una amenaza clandestina.
No todos cruzan para ganar ventaja,
algunos cruzan porque quedarse era la trampa.

[Pre-Chorus]
Y Maduro juega fuerte con la cámara encendida,
estado de excepción, redada dirigida.
Antes de las elecciones, fabrica un enemigo ideal,
la frontera convertida en escenario electoral.

[Chorus]
No me vendas dolor de patria,
si eliges qué dolor mirar.
Tres mil dentro de la casa,
¿quién los vino a reclamar?
No me vendas orgullo herido,
con la rabia en oferta especial.
Si la bandera tapa al vecino,
¿de qué sirve gritar nacional?

[Bridge]
Uribe, Pastrana, Samper, Santos,
cuántos años dejando el problema intacto.
El desplazado sin cámara, el migrante sin voz,
la misma deuda cambiando de color.

Santos y Holguín: hablen, negocien,
que ningún soldado la frontera desbroce.
Ni renuncia teatral ni desfile marcial,
la salida sensata también puede ser bilateral.

[Final Chorus]
No me vendas dolor de patria,
ni patriotero sentimental.
Mira al que pierde dentro de casa,
antes de posar para el altar.
No me vendas orgullo herido,
con titulares para incendiar.
Si la bandera tapa al vecino,
yo no la voy a levantar.

[Outro]
No todo grito es verdad,
no toda bandera sabe cuidar.
Primero la gente, después el papel,
la patria se demuestra cuando sabe responder.
\end{lyrics}

%! dir=wyomee

\songtitle{Entre mundos}
\tag*{2026}\tag{folk-rock}\tag{original-lyrics}\tag{spanish-lyrics}
\begin{notes}
Written by Carlos Thompson -- a lifelong habit of building alternate worlds as an escape from daily friction (traffic, bureaucracy, the gap between how things are and how they could be), and how generative AI became one more tool for exploring them. Retrieved from Wyomee Dank's Suno page (V6-MINI, 18 September 2026).
\end{notes}
\begin{style}
Rock with folk delivery, drum kit, ringing electric guitars and grounded bass, fast-paced driving groove, clean small-venue room sound with close-mic definition, mellow male vocals in a conversational folk cadence
\end{style}
\begin{lyrics}
[Estrofa 1]
Caminar la ciudad, conducir en tráfico de locos,
un sistema de transporte público entre tumultos o sofocado en el mismo trancón,
y saber que pudo haber sido diferente, administraciones con otros focos
o el escape a la fricción con mi propia imaginación.

[Estrofa 2]
Imaginarme cómo sería ese mundo ideal, imaginarme viviendo en él.
Planes elaborados desde lo grande en los sistemas hasta lo cotidiano.
Las diferencias en la fricción entre este mundo y el alterno en mi propia piel,
y pensar cómo pudo haber sido si pudiera hacerlo con mi propia mano.

[Puente]
La fantasía de poder sin dejar de ser quien soy.
No es por imaginarme ser el centro de atención.
Es más haber logrado ese otro mundo hoy
y vivir mi vida, pero así, sin fricción.

[Estribillo]
Entre mundos alternos y en mi propia vida
me quedo pensando más de lo prudente,
archivos y hojas que sobreviven la barrida,
lo que albergo aún en mi propia mente.

Entre mundos reales y los imaginarios
me quedo pensando sin poder evitarlo,
eventos mundanos y extraordinarios,
lo que sea si es que puedo pensarlo.

[Estrofa 3]
Leer un artículo en línea sobre cómo escribir y ahora quiero escribir.
Estudiar un nuevo lenguaje de programación y pensar cómo sería mi propio lenguaje.
Historia, biología, veo lo que hay y una parte de mí quiere vivir
en nuevas realidades, imaginar qué sería si, sea o no mi alterno viaje.

[Estrofa 4]
Llevo inventando mundos desde que me alcanza la memoria.
Jugar a ser quien no soy, pero también un yo diferente u otra ficción.
Otras lenguas, otros mundos, otros cuentos y también otra historia,
y buscando formas de poder materializar mi visión.

[Puente]
Llegan los modelos de lenguaje e IAs generativas,
nuevas formas de explorar el mundo,
asistentes y editores de mis narrativas.
Suficiente, si bien no tan profundo.

[Estribillo]
Entre mundos alternos y en mi propia vida
me quedo pensando más de lo prudente,
archivos y hojas que sobreviven la barrida,
lo que albergo aún en mi propia mente.

Entre mundos reales y los imaginarios
me quedo pensando sin poder evitarlo,
eventos mundanos y extraordinarios,
lo que sea si es que puedo pensarlo.

[Desglose]
Y sin darme cuenta tengo ya una novela
con todo: escenas, retratos, banda sonora.
He explorado mi mente, lo que teme y anhela,
y un universo alterno que no se demora.

He cantado a mis miedos, mis desvaríos,
a mis mundos alternos y mi fantasía,
a las mujeres reales que pensé mis amoríos,
a las mujeres imaginarias que dictó a la IA.

[Estribillo]
Entre mundos alternos y en mi propia vida
me quedo pensando más de lo prudente,
archivos y hojas que sobreviven la barrida,
lo que albergo aún en mi propia mente.

Entre mundos reales y los imaginarios
me quedo pensando sin poder evitarlo,
eventos mundanos y extraordinarios,
lo que sea si es que puedo pensarlo.

Entre mundos alternos,
entre mundos reales,
mi propia vida
y lo imaginario.
\end{lyrics}

\songtitle{Between Worlds}
\tag*{2026}\tag{folk-rock}\tag*{english-lyrics}\tag{entre-mundos}
\begin{notes}
English-language adaptation by Claude of \emph{Entre mundos}, by Carlos Thompson. Retrieved from Wyomee Dank's Suno page (V6-MINI, 18 September 2026).
\end{notes}
\begin{style}
Rock with mellow male folk-delivery vocals, close-mic and lightly doubled in choruses; driving drum kit, strummed and lightly overdriven guitars, melodic bass; fast-paced forward groove; clean small-venue room sound with restrained reverb and clear instrument separation.
\end{style}
\begin{lyrics}
[Verse 1]
Walking through the city, driving through a maddened tide,
a transit system swallowed whole or stuck in the same jam,
knowing it could've all gone different with a change in stride,
or the escape from friction into whatever world I am.

[Verse 2]
Imagining what shape that perfect world would wear,
imagining I'm living there, imagining I'm home.
Big designs from city systems down to what we wear,
the friction between this world and the one I roam.

[Pre-Chorus]
The fantasy of power, still the same soul underneath.
It's never been about the spotlight or the throne.
It's more like finally building that world beneath
and living out my life — but this time, friction-free, alone.

[Chorus]
Between the worlds I dream and the one I'm living in
I stay lost in thought, further than I should,
files and pages surviving every purge again,
what I still keep locked away for good.

Between the real world and the ones I've made
I stay lost in thought, I can't help but stray,
ordinary moments and the ones that don't fade,
whatever I can dream, come what may.

[Verse 3]
Read an article online on how to write, and now I want to write.
Study a new coding language, wonder what my own would sound like.
History, biology — I see what's here, but part of me tonight
wants new realities, wants to know what if, my alternate path or lifelike.

[Verse 4]
I've been building worlds since as far back as I recall.
Playing at not-me, at some other self, some other tale.
Other tongues, other worlds, other stories to enthrall,
searching for the way to make my vision real, not frail.

[Pre-Chorus]
Then come the language models, the generative minds,
new ways of walking through the world,
assistants, editors for the stories that I find.
Enough — even if it's not too deep unfurled.

[Chorus]
Between the worlds I dream and the one I'm living in
I stay lost in thought, further than I should,
files and pages surviving every purge again,
what I still keep locked away for good.

Between the real world and the ones I've made
I stay lost in thought, I can't help but stray,
ordinary moments and the ones that don't fade,
whatever I can dream, come what may.

[Breakdown]
And without noticing, I've got a whole novel grown —
scenes, portraits, a soundtrack of my own.
I've searched my mind, what it dreads and what it craves,
and a whole alternate universe that never stays.

I've sung to my fears, my wandering, restless mind,
to my other worlds and the fantasies I've spun,
to the real women I once called love, one kind,
to the imagined ones an AI dreamed, one by one.

[Final Chorus]
Between the worlds I dream and the one I'm living in
I stay lost in thought, further than I should,
files and pages surviving every purge again,
what I still keep locked away for good.

Between the real world and the ones I've made
I stay lost in thought, I can't help but stray,
ordinary moments and the ones that don't fade,
whatever I can dream, come what may.

Between the worlds I dream,
between the worlds that's real,
my own life
and the imaginary.
\end{lyrics}

\songtitle{Burden of Proof}
\tag*{2014}\tag{country-folk}\tag{original-source}
\begin{notes}
Lyrics by Suno based on a 2014 social-media reply by Carlos Thompson (Google+/YouTube), arguing burden of proof and the ad hominem fallacy in a comment-thread debate over religious claims.
\end{notes}
\begin{style}
Fast-paced country folk, layered lead and stacked harmony vocals, acoustic guitar strumming, banjo accents, upright bass, brushed snare and handclaps, conversational call-and-response phrasing, dry intimate verses opening into tight group harmonies, lively roots production, clear diction
\end{style}
\begin{lyrics}
[Verse 1]
Yes, what you are doing right now
Instead of facing the problem
You shift the topic around
And turn your aim upon the speaker

[Verse 2]
Whoever makes a claim
Has a burden to explain
You said that he, Jesus I suppose
Could be my very best friend
Then you offered quotes from a book
As the reason for the claim
But I've explained why I don't take
That book as an authority
And I've offered ample reasons
Though that's beyond the point
You made the claim, so you must show
Why your justification stands

[Chorus]
The burden stays with the one who claims
Not with the one who asks for proof
You can't demand I disprove your words
Then call my questioning the truth
The burden stays with the one who claims
Let the evidence carry the weight
If you ask me to take a step
You must show me why I should change

[Verse 3]
Now you ask if I understand
Burden of proof and ad hominem
As though I'm ignorant of the terms
And that settles the conversation
I'm not implying you are ignorant
You brought those terms into the room
I only think you're using them wrong
So here's how I understand them

[Verse 4]
When person A makes claim P
To person B across the way
The burden rests with person A
To justify what they say
It isn't B's duty to prove
That P is false or out of place
When A demands that B disprove it
A has shifted the burden's weight

[Bridge]
The negative is a claim as well
So which one carries proof?
The null hypothesis offers a path
To test what's most evident and true
One premise asks for less action
One asks me to become a friend
If you're the one who asks that step
The burden's yours to defend

[Verse 5]
Ad hominem means denying words
By pointing at the one who speaks
If you say the sky is blue
And I call you stupid, that's a cheap insult
But if I say the sky is not blue
Because you're stupid or blind
Your character does not change the sky
Or make the argument unkind

[Chorus]
The burden stays with the one who claims
Not with the one who asks for proof
You can't demand I disprove your words
Then call my question in the booth
The burden stays with the one who claims
Let the evidence carry the weight
If you ask me to take a step
You must show me why I should change

[Verse 6]
So far I haven't judged your character
Or Randal's, as a counterpoint
I've only said you may be mistaken
In the terms you chose to voice
I opened this thread to talk about
The skinniness of Healthyaddict
You and Randal came to promote something
But your strongest case is still that book

[Outro]
And you're offering it to people
Who don't take it as authority
We can agree to disagree
But that won't make your case for me
\end{lyrics}

\songtitle{The Question Left Between Us}
\tag*{2014}\tag{country-folk}\tag{original-source}\tag{burden-of-proof}
\begin{notes}
Lyrics by Suno based on the second part of the same 2014 social-media reply by Carlos Thompson (Google+/YouTube) set to \emph{Burden of Proof} -- respect, facetiousness, and misattributed statements in the same comment-thread debate, closing on an unrelated question about Ashley's vegan diet.
\end{notes}
\begin{style}
Fast-paced country folk, layered lead and harmony vocals, acoustic guitar strumming, banjo runs, brushed snare and handclaps, upright bass, conversational verses, punchy singalong refrains, warm live-room production, clear diction, dynamic vocal stacking, low variety
\end{style}
\begin{lyrics}
[Verse 1]
I have no interest in a pissing contest with you
Especially when your mind seems closed and your respect is few
I respect you and Randal, that's why I still reply
People deserve basic respect, and it can rise or die
You can earn it, you can lose it, that's the way I see
But ideas and beliefs don't get respect automatically

[Chorus]
Don't come preaching if you're not prepared to hear
Disagreement isn't hatred, it's a question standing clear
I'll listen for a reason, I'll change if you can show
But quoting from a book alone won't make authority grow

[Verse 2]
You say I'm being facetious, so I looked it up today
Funny how the dictionary changed the meaning on display
Then you claim I called all Christians something cruel and wrong
Find the place I ever said it; you won't search for long
I never made that statement, so please don't make it mine
Put the words back where they started, let the record draw the line

[Chorus]
Don't come preaching if you're not prepared to hear
Disagreement isn't hatred, it's a question standing clear
I'll listen for a reason, I'll change if you can show
But quoting from a book alone won't make authority grow

[Verse 3]
You say that's just opinion, but that isn't what I said
And an open mind is welcome, so here's the offer spread
Bring a solid reason with you, I'll consider what you bring
I've explained why the Bible isn't proof of everything
Give me one good reason I should take it as a guide
Not because it's being quoted, but because it can withstand the light

[Bridge]
You say you respect the atheist, understand the view
Then tell me what I believe and why I think the way I do
You don't call them names to win applause across the web
Good for you; that shallow habit stays exactly where you said

[Verse 4]
Seems like we found a point where both of us agree
So have a good day, and God bless, if that's your way to be
It's nighttime here, so have a morning, evening, or good night
And if you want the topic back, let's keep the question right
Do you think Ashley's vegan diet is healthy in the end?
That's the question left between us, friend

[Outro]
Do you think Ashley's vegan diet is healthy in the end?
That's the question left between us, friend
\end{lyrics}
